% TOP_PROBLEM: {"id":30007705,"problem_number":"LOCAL-30007705","title":"Continuous-treatment DiD","category_id":21,"category":{"id":21,"name":"economics","display_name":"Economics","description":"Open economic theory statements, published research agendas, and proposed scoped specifications; question provenance and historical priority are qualified per record.","slug":"economics","order_index":21,"created_at":"2026-10-08T00:00:00Z"},"status":"provenance_pending","external_url":null,"legacy_ids":[],"published":false,"statement":"Identify or bound selected-dose causal comparisons in a two-period panel using explicit restrictions beyond ordinary parallel trends.","economics":{"original_id":194,"field":"Econometrics, causal inference and measurement","kind":"identification","formalization_basis":"proposed_specification","source_status":"not_verified","priority_status":"not_established","priority_note":"No exact open-problem passage was verified. The supporting publication does not establish priority or unresolved status.","earliest_verified_reference":null,"current_reference":{"authors":["Brantly Callaway","Andrew Goodman-Bacon","Pedro H. C. Sant'Anna"],"title":"Difference-in-Differences with a Continuous Treatment","year":2024,"url":"https://www.nber.org/papers/w32117","doi":"10.3386/w32117","locator":"Primary indexed abstract","evidence_summary":"Establishes identified parameters and explains residual selection across doses; it supplies partial solutions rather than an explicit open declaration."},"formalization_note":"The cited primary work supports the subject or supplies solutions in particular settings, but no exact open-question passage for this specification was verified. The mathematical target and restrictions are proposed here, with no canonical-conjecture or priority claim."},"statusQualification":"Provenance pending: an explicit published statement of this exact open question was not verified. This is a catalogue-authored candidate specification, not a certified open conjecture. The cited primary work supports the subject or supplies solutions in particular settings, but no exact open-question passage for this specification was verified. The mathematical target and restrictions are proposed here, with no canonical-conjecture or priority claim.","statusReviewedAt":"2026-10-08"}
% FIRST_POSTED: 2026-10-08
% ENTRY_KIND: statement_only
% !TeX program = lualatex
\documentclass[11pt]{article}
\usepackage{fontspec}
\usepackage[a4paper,margin=25mm]{geometry}
\usepackage{amsmath,amssymb,mathtools}
\usepackage{xurl}
\usepackage[hidelinks]{hyperref}
\newcommand{\sref}[2]{\hyperlink{source:#1:#2}{[S#2]}}
\setlength{\emergencystretch}{3em}
\begin{document}
% problemId:30007705
\section*{Continuous-treatment DiD}\label{30007705}
\subsection{Definitions and mathematical statement}
Let units be observed before and after a policy, with baseline outcome \(Y_0\), nonnegative post-policy dose \(D\), and potential post-policy outcomes \(Y_1(d)\in[0,1]\) for feasible doses \(d\in[0,\bar d]\). Consistency gives observed \(Y_1=Y_1(D)\), and no anticipation makes \(Y_0\) unaffected by future dose. A conventional parallel-trends restriction compares \(E[Y_1(0)-Y_0\mid D=d]\) across dose groups; by itself it does not identify all comparisons between their treated potential outcomes. Fix a target dose interval \(B\) with positive probability and two doses \(d,d'\). Define
\[
\tau_B(d,d')=E[Y_1(d)-Y_1(d')\mid D\in B].
\]
Determine which additional selection restrictions, randomized dose information or bounded heterogeneity assumptions yield point identification or informative sharp bounds for this estimand from the observed panel law. For example, a declared Lipschitz constant \(L\) may restrict individual responses by \(|Y_1(d)-Y_1(d')|\leq L|d-d'|\), but its credibility requires external support. Distinguish dose-group average effects from causal derivatives, and state overlap and smoothness conditions whenever derivatives are targeted. The cited paper establishes several solutions under specified assumptions; this proposed restriction-and-bounds problem preserves residual selection concerns without describing continuous-treatment difference-in-differences as wholly unresolved.

\par\textbf{Formulation and scope.} The cited primary work supports the subject or supplies solutions in particular settings, but no exact open-question passage for this specification was verified. The mathematical target and restrictions are proposed here, with no canonical-conjecture or priority claim.

\par\textbf{Open-question provenance.} An explicit published open-question statement was not verified. Sources below are supporting literature and must not be treated as a first listing.

\par\textbf{Historical priority.} No exact open-problem passage was verified. The supporting publication does not establish priority or unresolved status.

\subsection{Short English statement}
Identify or bound selected-dose causal comparisons in a two-period panel using explicit restrictions beyond ordinary parallel trends.

\subsection{Sources}
\begin{itemize}\raggedright
\item[\textnormal{[S1]}]\hypertarget{source:30007705:1}{} Brantly Callaway, Andrew Goodman-Bacon, Pedro H. C. Sant'Anna. 2024. Difference-in-Differences with a Continuous Treatment. Primary indexed abstract.
\url{https://www.nber.org/papers/w32117}.
DOI: 10.3386/w32117.
{\small\textit{Supporting or current literature; see evidence qualification. Establishes identified parameters and explains residual selection across doses; it supplies partial solutions rather than an explicit open declaration.}}
\end{itemize}
\end{document}
